% DRENet 方法示意图 TikZ 版本
% 用于替代 ch3_drenet_overview.png

\begin{figure}[H]
  \centering
  \begin{tikzpicture}[
    box/.style={
      rectangle,
      rounded corners=2pt,
      draw=black,
      fill=white,
      minimum width=1.8cm,
      minimum height=0.7cm,
      text centered,
      font=\small
    },
    subbox/.style={
      rectangle,
      rounded corners=2pt,
      draw=black,
      fill=gray!10,
      minimum width=1.6cm,
      minimum height=0.6cm,
      text centered,
      font=\footnotesize
    },
    group/.style={
      rectangle,
      draw=black,
      dashed,
      rounded corners=3pt,
      inner sep=6pt
    },
    arrow/.style={
      -{Stealth[length=5pt, width=3.5pt]},
      thick,
      black
    },
    dashedarrow/.style={
      -{Stealth[length=5pt, width=3.5pt]},
      thick,
      black,
      dashed
    }
  ]
  
  % 顶部：主干网络流程
  \node[box] (input) at (0, 0) {Input image};
  \node[box] (backbone) at (2.5, 0) {Backbone};
  \node[box] (neck) at (5, 0) {Neck\\(CRMA)};
  \node[box] (head) at (7.5, 0) {Head\\(Detection)};
  \node[box] (output) at (10, 0) {Output};
  
  % 顶部箭头
  \draw[arrow] (input.east) -- (backbone.west);
  \draw[arrow] (backbone.east) -- (neck.west);
  \draw[arrow] (neck.east) -- (head.west);
  \draw[arrow] (head.east) -- (output.west);
  
  % 标注：shared features
  \node[font=\scriptsize] at (2.5, -0.9) {shared features};
  
  % 分支：退化重建增强
  \node[subbox] (degradation) at (5, -1.8) {Degradation};
  \node[subbox] (reconstruction) at (7, -1.8) {Reconstruction};
  
  % 从 backbone 到 degradation 的箭头
  \draw[arrow] (backbone.south) -- (degradation.west) node[midway, above, font=\tiny, sloped] {};
  
  % degradation 到 reconstruction
  \draw[arrow] (degradation.east) -- (reconstruction.west) node[midway, above, font=\scriptsize] {reconstruction loss};
  
  % reconstruction 回到 neck（特征引导，虚线）
  \draw[dashedarrow] (reconstruction.north) -- (neck.south) node[midway, right, font=\scriptsize, xshift=2pt] {feature guidance};
  
  % 分组框：Training Only
  \begin{scope}[on background layer]
    \node[group, fill=gray!5, fit=(degradation) (reconstruction), inner sep=10pt] (training) {};
    \node[font=\footnotesize, below] at (training.south) {Training only (not used in inference)};
  \end{scope}
  
  \end{tikzpicture}
  \caption{DRENet 方法示意图}
  \label{fig:ch3_drenet_overview}
\end{figure}
